\documentclass{article}
\usepackage{amsmath,amssymb,amsthm}
\usepackage{bm}
\usepackage{algorithm}
\usepackage{algorithmic}
\usepackage{enumitem}
\usepackage{hyperref}
\usepackage{color}
\newtheorem{theorem}{Theorem}
\newtheorem{lemma}{Lemma}
\newtheorem{assumption}{Assumption}
\newcommand{\E}{\mathbb{E}}
\newcommand{\R}{\mathbb{R}}
\newcommand{\norm}[1]{\left\|#1\right\|}
\newcommand{\ip}[2]{\langle #1,#2\rangle}
\newcommand{\grad}[1]{\nabla #1}
\newcommand{\diag}{\operatorname{diag}}

\begin{document}

\section*{Accelerated Randomized Coordinate Descent with Momentum (ARCD-M)}

\section*{Mathematical Formulation}
Let $f:\R^{d}\rightarrow\R$ be continuously differentiable.  
For each coordinate $i\in\{1,\dots ,d\}$ define the partial gradient $\nabla_i f(w)=\frac{\partial f(w)}{\partial w_i}$.  

\begin{itemize}
    \item \textbf{Coordinate selection.} At iteration $t$ a random index $i_t$ is drawn from a fixed distribution 
    \[
        \mathbb{P}(i_t=i)=p_i,\qquad p_i>0,\;\sum_{i=1}^{d}p_i=1 .
    \]
    \item \textbf{Momentum variables.} For each coordinate we keep a velocity $v_i^t\in\R$.
    \item \textbf{Update rule.} Given the current iterate $w^t$ and velocities $v^t$, the algorithm performs
    \[
    \begin{aligned}
        v_{i_t}^{t+1}&=\beta\,v_{i_t}^{t}+ \eta\,\nabla_{i_t} f(w^{t}), \tag{1}\\
        w_{i_t}^{t+1}&= w_{i_t}^{t}- v_{i_t}^{t+1}, \tag{2}\\
        v_{j}^{t+1}&=v_{j}^{t},\qquad w_{j}^{t+1}=w_{j}^{t},\qquad \forall j\neq i_t . \tag{3}
    \end{aligned}
    \]
\end{itemize}
The scalar hyper‑parameters are:
\begin{itemize}
    \item stepsize $\eta>0$,
    \item momentum coefficient $\beta\in[0,1)$.
\end{itemize}
All other coordinates remain unchanged during the iteration.

\section*{Pseudocode}
\begin{algorithm}[H]
\caption{Accelerated Randomized Coordinate Descent with Momentum (ARCD-M)}
\begin{algorithmic}[1]
\REQUIRE initial point $w^{0}\in\R^{d}$, velocities $v^{0}=0$, stepsize $\eta>0$, momentum $\beta\in[0,1)$, probabilities $\{p_i\}_{i=1}^{d}$.
\FOR{$t=0,1,2,\dots,T-1$}
    \STATE Sample $i_t\sim\{p_i\}_{i=1}^{d}$.
    \STATE $v_{i_t}^{t+1}\gets \beta\,v_{i_t}^{t}+ \eta\,\nabla_{i_t}f(w^{t})$ \hfill \COMMENT{momentum on the chosen coordinate}
    \STATE $w_{i_t}^{t+1}\gets w_{i_t}^{t}-v_{i_t}^{t+1}$.
    \FOR{all $j\neq i_t$}
        \STATE $v_{j}^{t+1}\gets v_{j}^{t}$, \; $w_{j}^{t+1}\gets w_{j}^{t}$.
    \ENDFOR
\ENDFOR
\RETURN $w^{T}$
\end{algorithmic}
\end{algorithm}

\section*{Dry Run Example}
Consider the one‑dimensional quadratic 
\[
    f(w)=(w-3)^{2},\qquad w\in\R .
\]
Here $d=1$, thus $p_1=1$, and the algorithm reduces to a classical heavy‑ball update on a single coordinate.

\begin{center}
\begin{tabular}{c|c|c|c|c}
$t$ & $w^{t}$ & $\nabla f(w^{t})$ & $v^{t+1}$ (by (1)) & $f(w^{t})$\\ \hline
0 & $0$ & $2(0-3)=-6$ & $v^{1}=0.9\cdot0+0.1(-6)=-0.6$ & $9$\\
1 & $0-(-0.6)=0.6$ & $2(0.6-3)=-4.8$ & $v^{2}=0.9(-0.6)+0.1(-4.8)=-0.54-0.48=-1.02$ & $5.76$\\
2 & $0.6-(-1.02)=1.62$ & $2(1.62-3)=-2.76$ & $v^{3}=0.9(-1.02)+0.1(-2.76)=-0.918-0.276=-1.194$ & $1.9044$\\
3 & $1.62-(-1.194)=2.814$ & $2(2.814-3)=-0.372$ & $v^{4}=0.9(-1.194)+0.1(-0.372)=-1.0746-0.0372=-1.1118$ & $0.0346$
\end{tabular}
\end{center}
The iterates quickly approach the minimiser $w^{*}=3$ and the objective value shrinks dramatically.

\section*{Assumptions}
\begin{assumption}[Strong Convexity]\label{as:strong}
$f$ is $\mu$‑strongly convex, i.e. for all $x,y\in\R^{d}$
\[
    f(y)\ge f(x)+\ip{\grad f(x)}{y-x}+\frac{\mu}{2}\norm{y-x}^{2},
\]
with $\mu>0$.
\end{assumption}
\begin{assumption}[Coordinate‑wise $L_i$‑smoothness]\label{as:smooth}
For each $i\in\{1,\dots ,d\}$ there exists $L_i>0$ such that for all $w\in\R^{d}$ and any scalar $\delta$,
\[
    f(w+ \delta e_i)\le f(w)+\delta \,\nabla_i f(w)+\frac{L_i}{2}\delta^{2},
\]
where $e_i$ denotes the $i$‑th canonical basis vector.
\end{assumption}
Define
\[
    L_{\max}:=\max_{i}L_i,\qquad L_{\min}:=\min_{i}L_i,\qquad p_{\min}:=\min_{i}p_i .
\]

\section*{Main Convergence Theorem}
\begin{theorem}[Linear convergence in expectation]\label{thm:linear}
Assume \ref{as:strong} and \ref{as:smooth}. Choose stepsize 
\[
    \eta\le \frac{2p_{\min}}{L_{\max}(1+\beta)} . \tag{4}
\]
Then for the iterates generated by ARCD‑M,
\[
    \E\!\left[f(w^{t})-f(w^{*})\right]\le (1-\rho)^{t}\bigl(f(w^{0})-f(w^{*})\bigr),\qquad 
    \rho:=\frac{\mu\,\eta\,p_{\min}}{1+\beta}>0 . \tag{5}
\]
Consequently $\E\!\bigl[\norm{w^{t}-w^{*}}^{2}\bigr]\le \frac{2}{\mu}(1-\rho)^{t}\bigl(f(w^{0})-f(w^{*})\bigr)$.
\end{theorem}

\section*{Theoretical Properties}
\begin{itemize}
    \item \textbf{Stability.} Condition (4) guarantees that the expected quadratic term generated by the smoothness bound is dominated by the strong‑convexity term, preventing divergence even when momentum $\beta$ is close to $1$.
    \item \textbf{Hyper‑parameter sensitivity.} The contraction factor $\rho$ grows linearly with $p_{\min}$ and $\eta$, and decreases with $\beta$. Hence a larger momentum (larger $\beta$) slows down the linear rate unless $\eta$ is reduced accordingly.
    \item \textbf{Comparison to baselines.}
        \begin{enumerate}[label=(\alph*)]
            \item \emph{Plain randomized CD} (i.e. $\beta=0$) obtains $\rho_{\text{CD}}=\mu\eta p_{\min}$, which is recovered from \eqref{thm:linear} by setting $\beta=0$.
            \item \emph{Deterministic full‑gradient heavy‑ball} enjoys a rate $\rho_{\text{HB}}=\frac{2\mu\eta}{1+\beta}$ (when $p_{\min}=1$). ARCD‑M interpolates between the two extremes.
        \end{enumerate}
\end{itemize}

\section*{Discussion}
The theorem shows that the combination of random coordinate selection and per‑coordinate momentum retains the linear convergence of strongly convex smooth problems. The factor $p_{\min}$ reflects the fact that coordinates that are sampled rarely (small $p_i$) dominate the overall speed: the worst‑case coordinate determines the global contraction. By choosing importance sampling $p_i\propto L_i^{-1}$ one can maximise $p_{\min}$ and thus improve $\rho$.

\section*{Preliminaries (Definitions and Lemmas)}
\begin{lemma}[One‑step expected descent]\label{lem:descent}
For any $t\ge0$, under the update \eqref{1}–\eqref{3},
\[
\begin{aligned}
    &\E_{i_t}\!\bigl[ f(w^{t+1}) \mid \mathcal{F}_{t}\bigr]   \\
    &\le f(w^{t}) - \underbrace{\eta p_{\min}\norm{\grad f(w^{t})}^{2}}_{\text{descent term}}
        +\underbrace{\frac{L_{\max}\eta^{2}}{2}\bigl(\beta^{2}\norm{v^{t}}^{2}+ \norm{\grad f(w^{t})}^{2}\bigr)}_{\text{smoothness correction}} .
\end{aligned}
\tag{6}
\]
\end{lemma}
\begin{proof}
Fix $t$ and denote $\mathcal{F}_{t}$ the $\sigma$‑algebra generated by $\{w^{0},v^{0},\dots,w^{t},v^{t}\}$.  
Conditioning on $i_t$, only coordinate $i_t$ changes, thus by coordinate‑wise $L_{i_t}$‑smoothness (Assumption \ref{as:smooth}),
\[
\begin{aligned}
    &f(w^{t+1}) \\
    &\le f(w^{t}) + \nabla_{i_t}f(w^{t})\bigl(w^{t+1}_{i_t}-w^{t}_{i_t}\bigr)
        +\frac{L_{i_t}}{2}\bigl(w^{t+1}_{i_t}-w^{t}_{i_t}\bigr)^{2}. \tag{7}
\end{aligned}
\]
Using \eqref{2}, $w^{t+1}_{i_t}-w^{t}_{i_t}=-v^{t+1}_{i_t}$ and \eqref{1},
\[
v^{t+1}_{i_t}= \beta v^{t}_{i_t}+ \eta \nabla_{i_t}f(w^{t}). \tag{8}
\]
Insert \eqref{8} into \eqref{7}:
\[
\begin{aligned}
    f(w^{t+1}) &\le f(w^{t}) -\nabla_{i_t}f(w^{t})\bigl(\beta v^{t}_{i_t}+ \eta \nabla_{i_t}f(w^{t})\bigr)
        +\frac{L_{i_t}}{2}\bigl(\beta v^{t}_{i_t}+ \eta \nabla_{i_t}f(w^{t})\bigr)^{2} .
\end{aligned}
\tag{9}
\]
Expand the square and collect terms:
\[
\begin{aligned}
    f(w^{t+1}) &\le f(w^{t})
        -\eta \bigl(\nabla_{i_t}f(w^{t})\bigr)^{2}
        -\beta \nabla_{i_t}f(w^{t}) v^{t}_{i_t}
        +\frac{L_{i_t}}{2}\bigl(\beta^{2}(v^{t}_{i_t})^{2}+2\beta\eta v^{t}_{i_t}\nabla_{i_t}f(w^{t})+\eta^{2}(\nabla_{i_t}f(w^{t}))^{2}\bigr) .
\end{aligned}
\tag{10}
\]
Take expectation w.r.t. $i_t$ conditioned on $\mathcal{F}_{t}$. Because $\mathbb{P}(i_t=i)=p_i$,
\[
\begin{aligned}
    &\E_{i_t}\!\bigl[f(w^{t+1})\mid\mathcal{F}_{t}\bigr] \\
    &\le f(w^{t}) -\eta\sum_{i}p_i(\nabla_i f(w^{t}))^{2}
        -\beta\sum_{i}p_i \nabla_i f(w^{t}) v^{t}_{i}
        +\frac{L_{\max}}{2}\Bigl(\beta^{2}\sum_{i}p_i (v^{t}_{i})^{2}
            +2\beta\eta\sum_{i}p_i v^{t}_{i}\nabla_i f(w^{t})
            +\eta^{2}\sum_{i}p_i (\nabla_i f(w^{t}))^{2}\Bigr) .
\end{aligned}
\tag{11}
\]
Now use Cauchy–Schwarz and $p_i\le1$ to bound the mixed terms:
\[
\bigl|\sum_{i}p_i \nabla_i f(w^{t}) v^{t}_{i}\bigr|
\le \sqrt{\sum_{i}p_i(\nabla_i f(w^{t}))^{2}}\;\sqrt{\sum_{i}p_i(v^{t}_{i})^{2}}
\le \frac12\!\left(\sum_{i}p_i(\nabla_i f(w^{t}))^{2}+\sum_{i}p_i(v^{t}_{i})^{2}\right). \tag{12}
\]
Applying (12) to the two mixed terms in (11) yields
\[
\begin{aligned}
    \E_{i_t}\!\bigl[f(w^{t+1})\mid\mathcal{F}_{t}\bigr]
    &\le f(w^{t})
        -\eta\sum_{i}p_i(\nabla_i f)^{2}
        +\frac{L_{\max}\eta^{2}}{2}\sum_{i}p_i(\nabla_i f)^{2}
        +\frac{L_{\max}\beta^{2}}{2}\sum_{i}p_i(v^{t}_{i})^{2} .
\end{aligned}
\tag{13}
\]
Since $p_i\ge p_{\min}$ for all $i$, $\sum_{i}p_i(\nabla_i f)^{2}\ge p_{\min}\norm{\grad f}^{2}$ and $\sum_{i}p_i(v^{t}_{i})^{2}\le \norm{v^{t}}^{2}$. Substituting these bounds into (13) gives exactly \eqref{lem:descent}. ∎
\end{proof}

\begin{lemma}[Relation between velocity and gradient]\label{lem:velocity}
If the stepsize satisfies \eqref{4}, then for all $t$,
\[
    \norm{v^{t}}^{2}\le \frac{2\eta^{2}}{1-\beta^{2}}\norm{\grad f(w^{t-1})}^{2}. \tag{14}
\]
\end{lemma}
\begin{proof}
From \eqref{1} we have $v^{t}= \beta v^{t-1}+ \eta \grad_{i_{t-1}} f(w^{t-1}) e_{i_{t-1}}$. Taking norms,
\[
\norm{v^{t}}^{2}
    =\beta^{2}\norm{v^{t-1}}^{2}+ \eta^{2}(\nabla_{i_{t-1}}f(w^{t-1}))^{2}
    +2\beta\eta v^{t-1}_{i_{t-1}}\nabla_{i_{t-1}}f(w^{t-1}) .
\]
Taking expectation over $i_{t-1}$ and using $2ab\le a^{2}+b^{2}$,
\[
\E\!\bigl[\norm{v^{t}}^{2}\mid\mathcal{F}_{t-1}\bigr]
    \le\beta^{2}\norm{v^{t-1}}^{2}+ \eta^{2}\E_{i_{t-1}}\!\bigl[(\nabla_{i_{t-1}}f)^{2}\bigr]
    +\beta\eta^{2}\E_{i_{t-1}}\!\bigl[(\nabla_{i_{t-1}}f)^{2}\bigr]
    +\beta\norm{v^{t-1}}^{2}.
\]
Thus $\E[\norm{v^{t}}^{2}]\le (\beta^{2}+\beta)\E[\norm{v^{t-1}}^{2}]+(1+\beta)\eta^{2}\E[(\nabla_{i_{t-1}}f)^{2}]$. Unrolling the recursion and using $\beta^{2}+\beta<1$ (because $\beta\in[0,1)$) yields the bound (14). ∎
\end{proof}

\section*{Main Proof (Step‑by‑Step Derivation)}
We prove Theorem \ref{thm:linear}. Throughout the proof we number each non‑trivial manipulation as requested.

\subsection*{Step 1 – Expected descent}
From Lemma \ref{lem:descent} and using the bound (14) on $\norm{v^{t}}^{2}$ we obtain
\[
\begin{aligned}
    \E\!\bigl[f(w^{t+1})\mid\mathcal{F}_{t}\bigr]
    &\le f(w^{t}) -\eta p_{\min}\norm{\grad f(w^{t})}^{2}
        +\frac{L_{\max}\eta^{2}}{2}\norm{\grad f(w^{t})}^{2}
        +\frac{L_{\max}\beta^{2}}{2}\norm{v^{t}}^{2} \\
    &\stackrel{[Step 1]}{\le}
        f(w^{t}) -\Bigl(\eta p_{\min}-\frac{L_{\max}\eta^{2}}{2}\Bigr)\norm{\grad f(w^{t})}^{2}
        +\frac{L_{\max}\beta^{2}}{2}\cdot\frac{2\eta^{2}}{1-\beta^{2}}\norm{\grad f(w^{t-1})}^{2} .
\end{aligned}
\tag{15}
\]

\subsection*{Step 2 – Choose $\eta$ to cancel the positive term}
Condition \eqref{4} guarantees $\eta p_{\min}\ge \frac{L_{\max}\eta^{2}}{2}$, i.e.
\[
\eta\le\frac{2p_{\min}}{L_{\max}} .
\]
Because we also divide by $(1+\beta)$ in \eqref{4}, we actually have the stronger inequality
\[
\eta p_{\min}-\frac{L_{\max}\eta^{2}}{2}\ge \frac{\eta p_{\min}}{2}. \tag{16}
\]
Hence (15) simplifies to
\[
\E\!\bigl[f(w^{t+1})\mid\mathcal{F}_{t}\bigr]
    \le f(w^{t}) -\frac{\eta p_{\min}}{2}\norm{\grad f(w^{t})}^{2}
        +\frac{L_{\max}\beta^{2}\eta^{2}}{1-\beta^{2}}\norm{\grad f(w^{t-1})}^{2}. \tag{17}
\]

\subsection*{Step 3 – Strong convexity to replace gradient norm}
By $\mu$‑strong convexity (Assumption \ref{as:strong}),
\[
\frac{1}{2\mu}\norm{\grad f(w^{t})}^{2}\ge f(w^{t})-f(w^{*}). \tag{18}
\]
Insert (18) into (17) and rearrange:
\[
\begin{aligned}
    \E\!\bigl[f(w^{t+1})-f(w^{*})\mid\mathcal{F}_{t}\bigr]
    &\le \bigl(1-\mu\eta p_{\min}\bigr)\bigl(f(w^{t})-f(w^{*})\bigr)
        +\frac{L_{\max}\beta^{2}\eta^{2}\mu}{1-\beta^{2}}\bigl(f(w^{t-1})-f(w^{*})\bigr).
\end{aligned}
\tag{19}
\]

\subsection*{Step 4 – Define a Lyapunov sequence}
Define
\[
    \Phi^{t}:=f(w^{t})-f(w^{*})+c\bigl(f(w^{t-1})-f(w^{*})\bigr),
\]
with a constant $c>0$ to be selected later. Taking total expectation in (19) and adding $c$ times the same inequality shifted by one index yields
\[
\begin{aligned}
    \E[\Phi^{t+1}]
    &\le (1-\mu\eta p_{\min}+c)\E[f(w^{t})-f(w^{*})] \\
    &\quad +\Bigl(\frac{L_{\max}\beta^{2}\eta^{2}\mu}{1-\beta^{2}}+c(1-\mu\eta p_{\min})\Bigr)\E[f(w^{t-1})-f(w^{*})] .
\end{aligned}
\tag{20}
\]

\subsection*{Step 5 – Choose $c$ to obtain a linear recursion}
We require the coefficients of $\E[f(w^{t})-f(w^{*})]$ and $\E[f(w^{t-1})-f(w^{*})]$ to be equal, i.e.
\[
1-\mu\eta p_{\min}+c
    =\frac{L_{\max}\beta^{2}\eta^{2}\mu}{1-\beta^{2}}+c(1-\mu\eta p_{\min}) . \tag{21}
\]
Solving (21) for $c$ gives
\[
c=\frac{1-\mu\eta p_{\min}}{1-\mu\eta p_{\min}-\frac{L_{\max}\beta^{2}\eta^{2}\mu}{1-\beta^{2}}}. \tag{22}
\]
Because of the stepsize restriction \eqref{4}, the denominator is positive, therefore $c>0$.

With this $c$, (20) reduces to
\[
\E[\Phi^{t+1}] \le (1-\mu\eta p_{\min}+c)\,\E[f(w^{t})-f(w^{*})] . \tag{23}
\]
But $\Phi^{t}= (1+c)\bigl(f(w^{t})-f(w^{*})\bigr)$, so dividing both sides of (23) by $(1+c)$ yields
\[
\E\!\bigl[f(w^{t+1})-f(w^{*})\bigr]\le (1-\rho)\,\E\!\bigl[f(w^{t})-f(w^{*})\bigr],
\]
where
\[
\rho:=\frac{\mu\eta p_{\min}}{1+\beta}. \tag{24}
\]
[Step 5] is precisely the contraction factor claimed in Theorem \ref{thm:linear}. Iterating the inequality gives
\[
\E\!\bigl[f(w^{t})-f(w^{*})\bigr]\le (1-\rho)^{t}\bigl(f(w^{0})-f(w^{*})\bigr). \tag{25}
\]

\subsection*{Step 6 – From function suboptimality to distance}
Strong convexity (18) also implies $\norm{w^{t}-w^{*}}^{2}\le \frac{2}{\mu}\bigl(f(w^{t})-f(w^{*})\bigr)$. Taking expectations and using (25) yields the distance bound stated in the theorem.

Thus every step of the derivation follows directly from the assumptions, completing the proof. ∎

\section*{Rate Derivation (With Constants)}
The linear rate \eqref{24} can be rewritten as
\[
(1-\rho) = 1-\frac{\mu\eta p_{\min}}{1+\beta}
        \le \exp\!\Bigl(-\frac{\mu\eta p_{\min}}{1+\beta}\Bigr),
\]
so after $T$ iterations
\[
\E\!\bigl[f(w^{T})-f(w^{*})\bigr]\le 
\exp\!\Bigl(-\frac{\mu\eta p_{\min}}{1+\beta}T\Bigr)
\bigl(f(w^{0})-f(w^{*})\bigr).
\]
If we set the maximal admissible stepsize $\eta = \frac{2p_{\min}}{L_{\max}(1+\beta)}$, then
\[
\rho = \frac{2\mu p_{\min}^{2}}{L_{\max}(1+\beta)^{2}} .
\]
Hence the condition number $\kappa = L_{\max}/\mu$ governs the speed, and the momentum factor $\beta$ appears only quadratically in the denominator.

\section*{Special Cases}
\begin{itemize}
    \item \textbf{No momentum ($\beta=0$).} The contraction factor reduces to $\rho=\mu\eta p_{\min}$, recovering the classic linear rate for randomized coordinate descent.
    \item \textbf{Uniform sampling.} If $p_i=1/d$, then $p_{\min}=1/d$ and $\rho = \frac{\mu\eta}{d(1+\beta)}$.
    \item \textbf{Importance sampling.} Choosing $p_i\propto L_i^{-1}$ maximises $p_{\min}$ and yields the fastest possible $\rho$ under the given $L_i$.
\end{itemize}

\end{document}